% How each method works: one short definition and one diagram per method family.
\tikzset{
  box/.style={draw, rounded corners=2pt, align=center, font=\footnotesize, minimum height=6mm, inner sep=3pt, fill=white},
  doc/.style={box, fill=gray!10},
  llm/.style={box, fill=red!10},
  model/.style={box, fill=blue!8},
  arr/.style={-{Stealth[length=2mm]}, semithick},
  lab/.style={font=\scriptsize\itshape, text=black!65}
}

\section{How each method works}\label{sec:methods}
This section defines each method in a few sentences and shows its data flow. $q$ is the query, $d$
a document, and every method ends with a ranked list of document ids.

\subsection{BM25 (lexical baseline)~\cite{kamphuis2020bm25,thakur2021beir}}
BM25 scores a document by the query terms it contains. Each matching term contributes its inverse
document frequency (rare terms count more), damped by term frequency (controlled by $k_1$) and
normalized by document length (controlled by $b$):
$\mathrm{score}(q,d)=\sum_{t\in q}\mathrm{idf}(t)\,\frac{tf_{t,d}(k_1+1)}{tf_{t,d}+k_1(1-b+b\,|d|/\overline{|d|})}$.
Text is tokenized and stemmed first, so ``running'' matches ``run'', but ``car'' never matches
``automobile''. An inverted index makes this very fast. BM25-BEIR is the same formula with BEIR's
parameters and with title and body scored separately and summed.
\begin{center}\begin{tikzpicture}[node distance=6mm]
\node[doc] (c) {Corpus};
\node[box, right=of c] (t) {Tokenize, stem,\\drop stopwords};
\node[box, right=of t] (i) {Inverted index\\term $\to$ docs, tf};
\node[box, right=of i] (s) {Sum BM25 term\\scores per doc};
\node[box, right=of s] (r) {Top 100};
\node[doc, above=4mm of t] (q) {Query $q$};
\draw[arr] (c)--(t); \draw[arr] (t)--(i); \draw[arr] (i)--(s); \draw[arr] (s)--(r);
\draw[arr] (q) -| (s);
\end{tikzpicture}\end{center}

\subsection{SPLADE (learned sparse retrieval)~\cite{formal2021splade,formal2021spladev2}}
SPLADE keeps BM25's ``bag of weighted terms'' form but lets a BERT model choose the terms and
weights. For every token position, the masked-language-model head predicts a score for each of
the 30,522 vocabulary terms. These are passed through $\log(1+\mathrm{ReLU}(\cdot))$ and max-pooled
over positions, giving one sparse vector per text. The vector also contains related terms that
never appear in the text (\emph{expansion}, e.g.\ ``car'' $\to$ ``vehicle''). Relevance is the dot
product of the query and document vectors, which an inverted index can compute.
\begin{center}\begin{tikzpicture}[node distance=6mm]
\node[doc] (d) {Text};
\node[model, right=of d] (b) {BERT + MLM head\\(score per vocab term)};
\node[box, right=of b] (p) {$\log(1{+}\mathrm{ReLU})$,\\max-pool};
\node[box, right=of p] (v) {Sparse vector\\\scriptsize car:2.1 vehicle:1.3 \dots};
\node[box, right=of v] (s) {$\mathbf{q}\cdot\mathbf{d}$};
\draw[arr] (d)--(b); \draw[arr] (b)--(p); \draw[arr] (p)--(v); \draw[arr] (v)--(s);
\node[lab, below=1mm of b] {same model for queries and documents};
\end{tikzpicture}\end{center}

\subsection{Dense bi-encoder (Dense = bge-small, Contriever)~\cite{karpukhin2020dpr,xiao2023cpack,izacard2022contriever}}
A transformer encodes the query and each document \emph{separately} into a single dense vector
(384 dimensions for bge-small, 768 for Contriever). Training places texts with the same meaning
close together, so a query and a relevant document match even if they share no words. Relevance is
cosine similarity (bge, normalized vectors) or a dot product (Contriever). Document vectors are
computed once, offline; at query time only the query is encoded, followed by a nearest-neighbour
search. bge-small is fine-tuned on large supervised relevance data; Contriever is trained without
labels (contrastive learning on text crops), which explains its weaker results.
\begin{center}\begin{tikzpicture}[node distance=6mm]
\node[doc] (d) {Documents};
\node[model, right=of d] (ed) {Encoder};
\node[box, right=of ed] (ix) {Vector store\\(offline)};
\node[doc, below=5mm of d] (q) {Query $q$};
\node[model, right=of q] (eq) {Encoder\\\scriptsize(same weights)};
\node[box, right=of eq] (qv) {Query vector};
\node[box, right=12mm of ix, yshift=-5.5mm] (s) {Cosine / dot\\nearest neighbours};
\node[box, right=of s] (r) {Top 100};
\draw[arr] (d)--(ed); \draw[arr] (ed)--(ix); \draw[arr] (q)--(eq); \draw[arr] (eq)--(qv);
\draw[arr] (ix)--(s); \draw[arr] (qv)--(s); \draw[arr] (s)--(r);
\end{tikzpicture}\end{center}

\subsection{ColBERTv2 with PLAID (late interaction)~\cite{khattab2020colbert,santhanam2022colbertv2,santhanam2022plaid}}
ColBERT keeps one vector \emph{per token} instead of one per text. A query is scored against a
document by MaxSim: each query token finds its most similar document token, and these maxima are
summed, $\mathrm{score}=\sum_{i\in q}\max_{j\in d}\mathbf{q}_i\cdot\mathbf{d}_j$. This captures
term-level matches with contextual meaning. Storing every token vector is expensive, so ColBERTv2
compresses each one to a cluster centroid plus a 2-bit residual. PLAID first uses the centroids to
pick a small candidate set cheaply, then computes exact MaxSim only for those candidates.
\begin{center}\begin{tikzpicture}[node distance=5mm]
\node[doc] (q) {Query tokens\\\scriptsize $q_1\dots q_{32}$};
\node[model, right=of q] (e) {BERT\\token vectors};
\node[box, right=of e] (c) {PLAID: centroid\\candidate pruning};
\node[box, right=of c] (m) {MaxSim per\\query token, sum};
\node[box, right=of m] (r) {Top 100};
\node[doc, below=5mm of c] (ix) {Compressed doc token vectors\\(centroid + 2-bit residual)};
\draw[arr] (q)--(e); \draw[arr] (e)--(c); \draw[arr] (c)--(m); \draw[arr] (m)--(r);
\draw[arr] (ix)--(c); \draw[arr] (ix) -| (m);
\end{tikzpicture}\end{center}

\subsection{Hybrid: Reciprocal Rank Fusion~\cite{cormack2009rrf}}
RRF merges several ranked lists using only the rank positions, so scores on different scales
(BM25 points vs.\ cosine similarities) never need calibrating. Each document receives
$\sum_{\ell}1/(60+\mathrm{rank}_\ell(d))$ over the lists that contain it. A document ranked highly
by either list, or moderately by both, rises to the top. Here the lists are the BM25 and Dense top 100.
\begin{center}\begin{tikzpicture}[node distance=6mm]
\node[box] (b) {BM25 top 100};
\node[box, below=3mm of b] (d) {Dense top 100};
\node[box, right=10mm of b, yshift=-4.5mm] (f) {$\sum 1/(60+\mathrm{rank})$\\per document};
\node[box, right=of f] (r) {Fused top 100};
\draw[arr] (b)--(f); \draw[arr] (d)--(f); \draw[arr] (f)--(r);
\end{tikzpicture}\end{center}

\subsection{Cross-encoder reranking (Rerank, BM25+CE)~\cite{thakur2021beir,wang2020minilm}}
A cross-encoder reads the query and one document \emph{together} in a single transformer pass
(``[CLS] query [SEP] document''), so every query word can attend to every document word. It outputs
one relevance score. This is more accurate than comparing separately computed vectors, but it must
run once per (query, document) pair at query time, so it is applied only to a shortlist from a
first-stage retriever: the Hybrid top 50 (Rerank) or the BM25-BEIR top 100 (BM25+CE).
\begin{center}\begin{tikzpicture}[node distance=6mm]
\node[box] (f) {First stage\\top $k$};
\node[model, right=10mm of f] (ce) {MiniLM cross-encoder\\$[q;\,d_i]\to s_i$};
\node[box, right=of ce] (s) {Sort by $s_i$};
\node[box, right=of s] (r) {Reranked top $k$\\+ rest unchanged};
\draw[arr] (f)--(ce); \draw[arr] (ce)--(s); \draw[arr] (s)--(r);
\end{tikzpicture}\end{center}

\subsection{HyDE*: Hypothetical Document Embeddings~\cite{gao2023hyde}}
\paragraph{Idea.} A dense retriever compares the query vector $f(q)$ with every document vector
$f(d)$ and returns the nearest. But queries and documents are different kinds of text: a query is a
few words (``Deleting Raptor reduces G-CSF levels.''), a document is a full abstract with its own
vocabulary. Their vectors live in different parts of the space, so the nearest documents to a
\emph{query} need not be the documents that \emph{answer} it. HyDE turns retrieval into
document--document similarity, which encoders are good at: an LLM \emph{writes} $N$ plausible answer
passages $\hat d_1,\dots,\hat d_N$ (``Please write a scientific paper passage to support/refute the
claim''), and the search looks for real documents that are close to those fake ones. The passages
may contain invented facts; that does not matter, because they are never shown to the user, only
encoded. The LLM never sees the corpus.

\begin{center}\begin{tikzpicture}[node distance=5mm]
\node[doc] (q) {Query $q$};
\node[llm, right=of q] (g) {LLM writes 8\\hypothetical passages};
\node[model, right=of g] (e) {Encoder\\$f(q),f(\hat d_1),\dots,f(\hat d_8)$};
\node[box, right=of e] (m) {Average\\vectors $\to \mathbf v$};
\node[box, right=of m] (s) {Nearest docs\\to $\mathbf v$ $\to$ top 100};
\draw[arr] (q)--(g); \draw[arr] (g)--(e); \draw[arr] (e)--(m); \draw[arr] (m)--(s);
\draw[arr] (q.south) -- ++(0,-4mm) -| node[lab, pos=0.25, below]{query also encoded} (e.south);
\end{tikzpicture}\end{center}
The search vector is the mean of $f(q)$ and the eight passage vectors, and the search itself is
the same nearest-neighbour search as Dense.

\noindent\textit{* HyDE is the method we plan to use for our project, CCDB Search (semantic search over the CFPB Consumer Complaint Database). How the averaged
vector finds the right documents, our measurements of it, and what they mean for CCDB Search are in
Section~\ref{sec:ccdb}.}

\subsection{RankGPT: listwise LLM reranking~\cite{sun2023rankgpt}}
RankGPT gives an LLM the query and 20 numbered passages and asks it to output a ranking such as
``[3] > [1] > [7] > \dots''. Because 100 passages do not fit well in one prompt, it uses a sliding
window: rerank positions 81--100, slide up by 10 to rerank 71--90, and so on to the top. Strong
documents can thus ``bubble up'' from the bottom of the list. That makes 9 LLM calls per query, which
must run in sequence. Unparseable or missing ids are appended in their original order.
\begin{center}\begin{tikzpicture}[node distance=5mm]
\node[box] (b) {BM25-BEIR\\top 100};
\node[llm, right=8mm of b] (w) {LLM ranks window\\of 20 passages};
\node[box, right=8mm of w] (p) {Parse ``[3] > [1] > \dots''\\reorder window};
\node[box, right=8mm of p] (r) {Reranked\\top 100};
\draw[arr] (b)--(w); \draw[arr] (w)--(p);
\draw[arr] (p.south) -- ++(0,-5mm) -| node[lab, pos=0.25, below]{slide window up by 10 (9 windows: 81--100 \dots 1--20)} (w.south);
\draw[arr] (p)--(r);
\end{tikzpicture}\end{center}

\subsection{LLM relevance judge (evaluation aid, not a retriever)}
Human relevance labels exist only for documents that were in the original judgment pool, and
unjudged documents count as non-relevant. Following Thomas et al.~\cite{thomas2024llmjudge}, an LLM (Qwen3.5) is shown the
query and one document and asked for a 0--2 relevance grade, after considering intent match (M) and
trustworthiness (T). We first measure its agreement with human labels, then use it only to label
the unjudged documents in each system's top 10. Human labels are never overwritten.
\begin{center}\begin{tikzpicture}[node distance=6mm]
\node[box] (u) {Unjudged (q, d)\\in any top 10};
\node[llm, right=of u] (j) {LLM judge\\0 / 1 / 2};
\node[box, right=of j] (a) {Human qrels\\+ LLM labels};
\node[box, right=of a] (e) {Re-score all\\systems};
\node[box, below=4mm of j] (v) {Validate vs.\ human\\labels ($\kappa$) first};
\draw[arr] (u)--(j); \draw[arr] (j)--(a); \draw[arr] (a)--(e); \draw[arr] (v)--(j);
\end{tikzpicture}\end{center}
